\documentclass{report}

\title{Whatsapp-Wrapped-CS108-Report}
\author{Kartik Bunas \& Vishal Gill}
\date{\today}

\begin{document}
\maketitle

\section{Division Of Work}
Kartik's Part:
\begin{itemize}
    \item Chat Generator
    \item HTML
    \item CSS
\end{itemize}
Vishal's Part:
\begin{itemize}
    \item Chat Analyzer
    \item JavaScript
\end{itemize}

\section{How To Use}
\begin{description}
	\item[Step 1] Navigate to the Whatsapp-Wrapped folder
	\item[Step 2] Start your python server by typing the following command "python3 -m http.server 8000"
	\item[Step 3] Open your browser and type the following address "http://localhost:8000/web/"
	\item[Step 4] Navigate the webpage .
\end{description}

\section{Directory Structure}

\begin{itemize}
    \item generator/
    \begin{itemize}
        \item chat.py
    \end{itemize}
    
    \item parser/
    \begin{itemize}
        \item analyze.py
    \end{itemize}
    
    \item report/
    \begin{itemize}
        \item main.pdf
        \item main.tex
        \item Makefile
        \item references.bib
    \end{itemize}
    
    \item web/
    \begin{itemize}
        \item index.html
        \item script.js
        \item style.css
    \end{itemize}
    
    \item chat.txt
    \item data.json
    \item README.md
    \item vocabulary.txt
\end{itemize}

\section{IMPORTANT NOTE}
The web/index.html includes script reference ,in line 8, to chart.js API which only works when we have internet connection . To make it work for offline mode , install the API and refer to it in the script tag instead to the url

\chapter{By Kartik}

\section{Libraries Used}

\subsection{Chat Generator}
\begin{description}
  \item[datetime] Used to describe the start and end times, time increment, and format the time string.
  \item[numpy] Used to randomize the sender, the gaps between messages, and the message wording.
  \item[pathlib] Used to choose the location to save chat.txt file generated.
  \item[sys] Used to read command line arguments.
\end{description}

\section{Features Implemented}
Our Wrapped Project allows us to generate and analyze chat messages. We can view this as a slideshow in which there are some custom awards, charts, etc. for the whole group and at the end we have the option to analyze each member's activity.

\section{Hurdles Faced}
\subsection{Chat Generator}
\begin{itemize}
    \item I didn't get how do I save the generated text into chat.txt which was not in the same folder as the generator . This is where pathlib came in handy. I looked it up on \cite{pythondocs}.
    \item I didn't know what all numpy distributions I had to use for the randomness in time manipulation and forming messages. I looked this up online too.
\end{itemize}
\subsection{HTML \& CSS}
There wasn't any explicit major issues. There were some concepts I didn't know about flexbox, transitions, padding/margins, etc . I used \cite{gemini} and \cite{w3schools} for this and also for looking up the colours used in spotify.The browser also refreshed on each slide so I looked up the solution for it which was hidden view.
\subsection{LaTeX}
Used \cite{gemini} to make references file and for help in pdf compiling using makefile.

\chapter{By Vishal}

\section{Libraries Used}

\subsection{Chat Analyzer}
\begin{description}
  \item[sys] Used to read the chat.txt file path from the command line arguments.
  \item[datetime] Used to parse timestamp strings into real datetime objects and compute time differences between messages.
  \item[re] Used to detect and extract emojis from messages using unicode ranges \cite{pythondocs}.
  \item[numpy] Used to compute the median response times for each person.
  \item[json] Used to write all computed statistics into data.json.
  \item[pathlib] Used to locate and save the data.json file relative to the script location.
\end{description}

\subsection{JavaScript}
\begin{description}
  \item[Chart.js] Used to render bar charts for message counts, conversation starters, response times, and night owl activity \cite{chartjs}.
\end{description}

\section{Features Implemented}

\subsection{Chat Analyzer}
The analyzer reads chat.txt line by line and parses each line into three components — a datetime object, a sender name, and a message string. It then computes the following statistics:

\begin{itemize}
    \item \textbf{Total Messages:} Count of messages sent by each person.
    \item \textbf{Total Words:} Total number of words typed by each person.
    \item \textbf{Night Owl:} Person who sent the most messages between 12am and 4am.
    \item \textbf{Ghost:} Person whose messages received the fewest replies. A reply is defined as a message sent by a different person within 30 minutes.
    \item \textbf{Conversation Starter:} Person who most frequently sent a message after a silence of more than 60 minutes.
    \item \textbf{Most Used Emoji:} Top 3 emojis per person, detected using unicode range pattern matching.
    \item \textbf{Busiest Day:} The single date with the highest number of messages.
    \item \textbf{Longest Silence:} The largest time gap between two consecutive messages, reported in hours.
    \item \textbf{Average Response Time:} Median time in minutes before someone replied to each person's messages.
    \item \textbf{Hype Person:} Person with the fastest median reply time to others.
    \item \textbf{Custom Stat - Max Consecutive Streak:} The longest streak of back-to-back messages sent by one person without anyone else replying. This captures the "ghoster" or "ranter" personality unique to this chat.
\end{itemize}

\subsection{JavaScript}
The script.js file loads data.json using fetch() and populates all 9 slides dynamically \cite{mdn}. It handles slide navigation, draws all charts using Chart.js, builds the emoji grid, and implements the per-person profile selector where the user can click any member's name to view their individual statistics.

\section{Custom Statistic}
The custom statistic I chose is \textbf{Max Consecutive Streak} — the longest number of messages a single person sends in a row without anyone else replying. This metric is interesting because it captures a personality pattern unique to our simulated group chat. A high streak indicates someone who "rants" or "ghosts" the group by flooding messages without waiting for responses. A generic tool would only count total messages, but this stat reveals \textit{how} someone messages, not just \textit{how much}.

\section{Definition Choices in the Parser}

\begin{itemize}
    \item \textbf{Direct Reply:} A message sent by a \textit{different} person within 30 minutes of the previous message. 30 minutes was chosen as it reflects a realistic window for active conversation.
    \item \textbf{Long Silence:} A gap of more than 60 minutes between two consecutive messages. Below 60 minutes feels like a natural pause, above it feels like the conversation has ended.
    \item \textbf{Busiest Day:} The day with the highest raw message count. This measures volume, not spread — a day with 100 messages in one hour counts more than a day with 80 messages spread across 12 hours.
\end{itemize}

\section{Hurdles Faced}

\begin{itemize}
    \item Parsing each line of chat.txt was tricky at first. The format had to be split twice — first by \texttt{" - "} to separate the timestamp, then by \texttt{": "} to separate the sender from the message. I had to use the \texttt{maxsplit=1} argument to avoid breaking messages that contained colons.
    \item Computing the ghost stat required careful thinking about what counts as a reply. I had to loop through consecutive message pairs and check both the sender and the time gap simultaneously.
    \item Detecting emojis was not straightforward. I used the \texttt{re} module with unicode ranges since standard string methods don't handle them well \cite{pythondocs}.
    \item In JavaScript, the profile chart would crash if you viewed one profile, went back, and clicked another — because Chart.js does not allow drawing on a canvas that already has a chart. I fixed this by destroying the old chart instance before drawing a new one \cite{mdn}.
\end{itemize}

\section{What I Would Do Differently With 10 More Hours}

\begin{itemize}
    \item Add hourly breakdown data to data.json so the activity chart shows a proper 24-hour heatmap instead of a per-person bar chart.
    \item Make the per-person profile slide show the person's top emojis visually instead of just numbers.
    \item Add animations between slides to make the Wrapped experience feel more like Spotify's version.
    \item Write more robust error handling in the parser for edge cases like missing timestamps or corrupted lines.
\end{itemize}


\bibliographystyle{plain}
\bibliography{references}
\end{document}
